\subsubsection{Spatial Filters}\label{sec:spatial-filters}

In the previous section we introduced the idea of a local neighbourhood. For every pixel $\left(r,c\right)$, we considered a small window $U_{r,c}$ containing only the nearby pixels. The natural next question is: \emph{\textbf{what should we do with those pixels?}} The answer is that we can apply a \textbf{spatial filter}, which is a function that receives the neighbourhood and computes the output pixel.

\highspace
The filter $T_U$ receives all the pixels inside the neighbourhood $U_{r,c}$ and returns one new pixel value. Formally, we can write this as:
\begin{equation}
    T_U: U_{r,c} \to G\left(r,c\right)
\end{equation}
Where $G\left(r,c\right)$ is the output pixel value at position $\left(r,c\right)$. The transformation could be the average of the pixels, the median, or any other function that takes a set of pixel values and returns a single value. The important point is \textbf{not yet how we compute the value}, but \textbf{that the computation uses only the local neighbourhood}.

\highspace
\begin{flushleft}
    \textcolor{Green3}{\faIcon{book} \textbf{Space-Invariant Transformations}}
\end{flushleft}
The spatial filter $T_U$ is said to be \textbf{space-invariant} if it does not depend on the position of the pixel $\left(r,c\right)$. In other words, \hl{the same function is applied everywhere in the image}. This property is important because it allows us to apply the same filter to different parts of the image without having to redefine it for each location.

\highspace
For example, imagine an edge detector. If the filter detects vertical edges, we want it to detect them at any location in the image, not just at a specific row or column. This is what makes the filter space-invariant. In simple terms, with this property, the \textbf{input values change} (the pixel values in the neighbourhood), but the \textbf{function itself does not change} (the way we compute the output pixel value).

\highspace
\textcolor{Green3}{\faIcon{tools} \textbf{General Pipeline.}} The complete processing pipeline is now:
\begin{enumerate}
    \item Input image $I$.
    \item Select a neighbourhood $U_{r,c}$ for each pixel $\left(r,c\right)$.
    \item Apply a spatial filter $T_U$ to the neighbourhood $U_{r,c}$ to compute the output pixel value $G\left(r,c\right)$.
    \item Move the sliding window to the next pixel.
    \item Repeat the process until all pixels have been processed.
    \item Output image $G$.
\end{enumerate}
Every pixel of the output image is obtained by repeating exactly this procedure.